\chapter{Prédiction et market making sur les marchés sportifs}\label{ch:hf:prediction-and-sports-market-making}

Un but est marqué, et l'image de la télévision arrive plusieurs secondes après celle du stade. La bourse de paris retient chaque ordre en cours de match pendant quelques
secondes avant de l'appairer, afin que les supporters dans les tribunes ne gagnent pas à tous les coups. Dans le modèle de ce chapitre, un teneur de marché sur un marché du vainqueur
dont le flux de données signale les buts trois secondes après qu'ils surviennent perd 236~\$ par but face aux courtsiders présents au stade lorsque la bourse appaire immédiatement~;
un \omterm{def:hf:prediction-and-sports-market-making:delay}{délai de pari} de trois secondes réduit cette perte à 66~\$, un délai de six secondes à 1~\$, et son revenu de fourchette de 90~\$ par match devient un profit.

\section{Modèles de prix pour les événements et les matchs}

Les marchés de prédiction et les bourses de paris (livre~3, chapitre~27) négocient des contrats qui versent un montant si un événement se produit~; leurs prix sont
des probabilités. Un match de football se modélise assez bien, pour un teneur de marché, par des buts de Poisson indépendants pour chaque équipe (livre~4, chapitre~6)~: avec
des intensités de 1.5 et 1.1 but par match, l'équipe à domicile l'emporte avec une probabilité de 0.464 au coup d'envoi. En cours de jeu, la probabilité est recalculée à partir du
score et du temps restant~: une avance d'un but à la mi-temps donne 0.760. Chaque but est un saut du \omterm{def:hf:fair-value:fair}{prix juste}, et tout le problème du teneur de marché se joue
dans les secondes qui suivent (\cref{fig:hf:prediction-and-sports-market-making:match}).

\begin{omfigure}
\begin{tikzpicture}
\begin{axis}[width=11cm,height=5cm,xlabel={\small minute de jeu},ylabel={\small probabilité de victoire à domicile},xmin=0,xmax=90,ymin=0,ymax=1.05,
  tick label style={font=\footnotesize},grid=major,grid style={gray!25},table/col sep=comma]
\addplot[omThm,thick,const plot] table[x=minute,y=p]{figdata/hft/26-prediction-and-sports-market-making/match.csv};
\draw[omDef,dashed] (axis cs:0.47,0) -- (axis cs:0.47,1.05);
\draw[omDef,dashed] (axis cs:73.9,0) -- (axis cs:73.9,1.05);
\draw[omDef,dashed] (axis cs:78.6,0) -- (axis cs:78.6,1.05);
\end{axis}
\end{tikzpicture}

\omcaption{La probabilité juste de victoire à domicile, minute par minute, au fil d'un match simulé avec des buts de Poisson à 1.5 et 1.1 par match~; les pointillés
marquent les buts de l'équipe à domicile aux 1\up{re}, 74\up{e} et 79\up{e} minutes. La probabilité dérive au fil du temps et saute à chaque but.
Données~: \texttt{hf\_sports.match\_path}.}\label{fig:hf:prediction-and-sports-market-making:match}
\end{omfigure}

Croxson et Reade ont utilisé exactement cette information nette, un but, pour tester l'efficience des prix des paris sur des données à haute fréquence~: selon leurs
résultats, les prix se mettent à jour rapidement et intégralement. La rapidité est précisément l'enjeu~: il faut arriver le premier.

\section{Faire un marché sur les bourses et les plateformes de prédiction}

Un teneur de marché sur une bourse cote les deux côtés d'un contrat~: il affiche la victoire à domicile à la probabilité juste plus une demi-fourchette, et la reprend à la
probabilité juste moins un point. À 0.46 avec une demi-fourchette d'un point, un parieur obtient une cote décimale de 2.13 sur la victoire à domicile et de 1.82 sur l'autre issue, un
book avec un overround de 2~\% (\texttt{firm.odds} du livre~3). Sur une plateforme de prédiction organisée en carnet d'ordres à cours limité centralisé, la même cotation prend la forme d'un cours acheteur
et d'un cours vendeur sur un contrat dont le prix est compris entre 0 et 1. Les parieurs occasionnels paient la fourchette~; la mise répétée deux fois par minute à 50 de ce chapitre donne au
teneur de marché 90~\$ par match.

\begin{dated}{2026-09}{Où ces contrats se négocient, et le délai}\label{dat:hf:prediction-and-sports-market-making:venues}
Le service d'assistance aux développeurs d'une bourse de paris explique que les marchés en cours de match imposent un délai de placement des paris, variant de 1 à 12 secondes, afin de permettre aux clients
d'annuler les ordres non appariés lorsque les conditions de marché changent. Aux États-Unis, un tribunal fédéral de district a jugé le 12 septembre 2024 que la CFTC avait
excédé ses pouvoirs en bloquant les contrats d'élection d'un designated contract market~; l'appel de la CFTC a été retiré volontairement le 7 mai 2025.
En janvier 2022, la CFTC a réglé des poursuites contre une plateforme de marché de prédiction crypto pour avoir proposé des options binaires fondées sur des événements sans enregistrement, avec une
pénalité de 1.4~million de dollars (livre~3, chapitre~27).
\end{dated}

\section{Latence en cours de match et délai de pari}

\begin{definition}[Courtsiding]\label{def:hf:prediction-and-sports-market-making:courtsiding}
Le \emph{courtsiding}\index{courtsiding} consiste à trader sur un événement depuis l'intérieur du lieu où il se déroule, plus vite que le flux de données officiel et la
retransmission télévisée ne le transmettent, afin de saisir des cotations qui ne le reflètent pas encore.
\end{definition}

\begin{definition}[Délai de pari]\label{def:hf:prediction-and-sports-market-making:delay}
Un \emph{délai de pari}\index{délai de pari} (bet delay) est le temps pendant lequel une bourse de paris retient un ordre en cours de match avant de pouvoir l'appairer, afin que les teneurs de marché puissent annuler
leurs cotations après un événement, avant que les ordres passés sur la connaissance de cet événement ne les atteignent~; c'est un \omterm{def:hf:latency-arbitrage-and-its-defence:bump}{ralentisseur asymétrique} (chapitre~9) pour le sport.
\end{definition}

\begin{omfigure}
\begin{tikzpicture}[font=\footnotesize,>=stealth]
\draw[->] (0,0) -- (12,0) node[right] {secondes};
\foreach \x/\l in {0/0,2/1,4/2,6/3,8/4,10/5}{\draw (\x,0.1) -- (\x,-0.1) node[below] {\l};}
\draw[omDef,very thick] (0,0.2) -- (0,1.9) node[above] {but};
\draw[omThm,thick,->] (1,0.5) -- (1,1.3) node[above,align=center] {courtsider\\envoie};
\draw[gray,thick,->] (1,1.0) -- (7,1.0) node[midway,above] {délai de pari (3 s)};
\draw[omThm,very thick] (7,0.2) -- (7,1.6) node[above] {ordre peut être apparié};
\draw[black,thick] (6,0.2) -- (6,0.9);
\node[align=center,anchor=north] at (6,-0.5) {flux du teneur\\signale~; annule};
\end{tikzpicture}

\omcaption{La course après un but~: un courtsider présent au stade envoie un ordre après une demi-seconde~; la bourse retient l'ordre pendant le \omterm{def:hf:prediction-and-sports-market-making:delay}{délai de pari}~; le
flux du teneur de marché signale le but après environ trois secondes, et celui-ci annule aussitôt. Avec un \omterm{def:hf:prediction-and-sports-market-making:delay}{délai de pari} de trois secondes, l'annulation du teneur, à 3 secondes,
arrive avant que l'ordre puisse être apparié, à 3.5~; avec un délai de deux secondes, ce n'est pas le cas. Les latences varient autour de ces médianes. Source~: ce
chapitre.}\label{fig:hf:prediction-and-sports-market-making:race}
\end{omfigure}

La course du chapitre (\cref{fig:hf:prediction-and-sports-market-making:race}, \cref{lst:hf:prediction-and-sports-market-making:trade})~: à chaque but,
le courtsider agit après une médiane d'une demi-seconde, et le flux du teneur après une médiane de trois secondes, toutes deux lognormales avec un écart-type logarithmique de 0.3~;
l'ordre du courtsider, retenu pendant le délai, prend les 1\,000 unités du teneur à la cotation périmée s'il arrive avant l'annulation. Le teneur
perd la mise multipliée par le saut de probabilité, diminuée de sa demi-fourchette.

\begin{omfigure}
\begin{tikzpicture}
\begin{axis}[width=11cm,height=5.2cm,axis y line*=right,axis x line=none,xmin=0,xmax=10,ymin=0,ymax=105,
  ylabel={\small buts saisis (\%)},tick label style={font=\footnotesize},table/col sep=comma]
\addplot[omDef,thick,dashed,mark=square*,mark size=1.3pt] table[x=delay,y=sniped]{figdata/hft/26-prediction-and-sports-market-making/delay.csv};
\label{plot:hf:sniped}
\end{axis}
\begin{axis}[width=11cm,height=5.2cm,axis y line*=left,xmin=0,xmax=10,xlabel={\small délai de pari (secondes)},ylabel={\small perte par but (\$)},
  ymin=0,ymax=260,tick label style={font=\footnotesize},grid=major,grid style={gray!25},table/col sep=comma,
  legend style={font=\scriptsize,at={(0.97,0.55)},anchor=east}]
\addplot[omThm,thick,mark=*,mark size=1.4pt] table[x=delay,y=loss]{figdata/hft/26-prediction-and-sports-market-making/delay.csv};
\addlegendentry{perte par but (gauche)}
\addlegendimage{/pgfplots/refstyle=plot:hf:sniped}\addlegendentry{part des buts saisis (droite)}
\draw[gray,dashed] (axis cs:3,0) -- (axis cs:3,260);
\node[font=\scriptsize,anchor=south west] at (axis cs:3.1,215) {latence médiane du flux du teneur};
\end{axis}
\end{tikzpicture}

\omcaption{La perte du teneur de marché par but face aux courtsiders et la part des buts sur lesquels ceux-ci atteignent sa cotation périmée, en fonction du \omterm{def:hf:prediction-and-sports-market-making:delay}{délai de pari} de la bourse~;
2\,000 matchs simulés (5\,241 buts), courtsiders à une demi-seconde et teneur à trois secondes après chaque but (lognormal). Données~:
\texttt{hf\_sports.by\_delay}.}\label{fig:hf:prediction-and-sports-market-making:delay}
\end{omfigure}

Sans délai, chaque but est saisi par les courtsiders et le teneur perd 236~\$ par but, soit 528~\$ par match net de son revenu de fourchette
(\cref{fig:hf:prediction-and-sports-market-making:delay}). Les délais plus courts que la propre latence du teneur aident peu~: à deux secondes, 72~\% des buts
sont encore saisis. À trois secondes, 29~\% et 66~\$~; à quatre, 8~\% et 20~\$, et le marché rapporte~; à six, 1~\$ par but, soit moins de 1~\% de la perte sans délai.
Le délai doit dépasser la latence des données du teneur de marché diminuée de celle du courtsider, avec une marge pour la fourchette des deux parties.

\section{Limites et sélection des contreparties}

\begin{definition}[Limitation des mises]\label{def:hf:prediction-and-sports-market-making:limiting}
La \emph{limitation des mises}\index{limitation des mises} (stake limiting) est la restriction, par un bookmaker ou un teneur de marché, des mises qu'il accepte d'un compte dont les paris se sont révélés informés,
mesurée par le mark-out de ses paris passés, et qui laisse ses cotes ouvertes aux autres.
\end{definition}

Un bookmaker voit le compte sur chaque pari et peut le limiter~; un teneur de marché d'une bourse ne voit que des ordres anonymes, mais une plateforme peut lui offrir des outils
(délais minimaux par marché, plafonds par compte). Dans le modèle, plafonner la mise de chaque courtsider à 100 unités au lieu des 1\,000 cotées réduit la perte
par but de 236~\$ à 24~\$, et à 50 unités à 12~\$. Le \omterm{def:hf:prediction-and-sports-market-making:courtsiding}{courtsiding} se situe lui-même dans une zone grise juridique~: en 2014, le parquet de Victoria a abandonné une
accusation de comportement susceptible de fausser l'issue d'un pari contre un homme accusé d'avoir transmis des scores en direct de l'Open d'Australie à une agence de paris,
faute de perspective raisonnable de condamnation.

\section{Fiches de stratégie}

\begin{strategyfile}{Market making sur marché de prédiction, sur un carnet d'ordres à cours limité central}\label{strat:hf:prediction-and-sports-market-making:clob}
\sfield{Qui paie, et pourquoi} Les traders qui expriment des vues sur des événements, et les hedgers du risque d'événement, qui paient la fourchette.
\sfield{Instruments et plateformes} Contrats d'événement sur des designated contract markets (DCM) et plateformes de prédiction crypto.
\sfield{Signal} Une probabilité issue d'un modèle (sondages, modèles, marchés liés) et le flux du carnet d'ordres.
\sfield{Dimensionnement et exécution} Coter des deux côtés autour du modèle, introduire un skew selon l'inventaire, plafonner l'exposition par événement, élargir la fourchette à l'approche des annonces de résolution.
\sfield{Coûts} Frais~; capital bloqué jusqu'à la résolution~; risque lié à la source de résolution.
\sfield{Comment elle échoue} Flux informé à l'approche des annonces, et réglementation~: les décisions de justice et le règlement-livraison de l'encadré daté montrent à quelle vitesse l'ensemble des opérations autorisées évolue.
\sfield{Horizon, capacité, infrastructure} De quelques heures à quelques mois jusqu'à la résolution.
\sfield{Backtest en toute honnêteté} Le carnet de la plateforme à chaque instant, avec ses frais~; la résolution telle qu'elle a eu lieu.
\sfield{Sources} le livre~3, chapitre~27~; l'encadré daté.
\end{strategyfile}

\begin{strategyfile}{Négociation en direct sur bourse de paris avec délai de pari}\label{strat:hf:prediction-and-sports-market-making:inplay}
\sfield{Qui paie, et pourquoi} Les parieurs récréatifs en direct, via la fourchette~; le délai protège le teneur de marché des courtsiders.
\sfield{Instruments et plateformes} Marchés de match-odds et autres marchés en direct sur les bourses de paris.
\sfield{Signal} Un modèle en direct (buts de Poisson pour le football) mis à jour à partir d'un flux de données officiel.
\sfield{Dimensionnement et exécution} Coter autour du modèle~; annuler à chaque événement issu du flux~; dimensionner selon la marge du délai sur la latence du flux.
\sfield{Coûts} Commission de la bourse sur les gains nets~; pertes lorsque le flux est plus lent que le délai.
\sfield{Comment elle échoue} Délais plus courts que la latence du flux~: 236~\$ par but sans délai, 66~\$ à trois secondes~; données plus rapides pour les courtsiders.
\sfield{Horizon, capacité, infrastructure} Secondes~; le flux officiel le plus rapide disponible.
\sfield{Backtest en toute honnêteté} Horodatages des événements relevés sur le terrain, ceux du flux et ceux de la bourse, et délai par marché.
\sfield{Sources} Croxson et Reade (2014)~; le support aux développeurs de la bourse (encadré daté)~; le présent chapitre.
\end{strategyfile}

\begin{strategyfile}{Arbitrage d'événements inter-plateformes}\label{strat:hf:prediction-and-sports-market-making:cross}
\sfield{Qui paie, et pourquoi} Les traders de la plateforme dont le prix accuse un retard, ou les parieurs auprès des bookmakers dont l'overround laisse de la marge.
\sfield{Instruments et plateformes} Le même événement sur des bourses, chez des bookmakers et sur des plateformes de prédiction.
\sfield{Signal} Probabilités implicites entre plateformes, nettes des commissions et des overrounds (\texttt{firm.odds} du livre~3).
\sfield{Dimensionnement et exécution} Back sur une plateforme et lay sur une autre lorsque les cotes combinées garantissent un profit~; petites mises là où la limitation se fait sentir.
\sfield{Coûts} Commissions, \omterm{def:hf:prediction-and-sports-market-making:limiting}{limitations des mises}, capital réparti sur plusieurs plateformes, règles de règlement-livraison différentes.
\sfield{Comment elle échoue} \omterm{def:hf:prediction-and-sports-market-making:limiting}{Limitation des mises} des comptes gagnants, et différences dans la façon dont les plateformes règlent un même événement.
\sfield{Horizon, capacité, infrastructure} De quelques minutes jusqu'à l'événement.
\sfield{Backtest en toute honnêteté} Cotes accessibles au compte à l'instant considéré, commissions et règles de chaque plateforme.
\sfield{Sources} le livre~3, chapitre~27.
\end{strategyfile}

\section{Tutoriel~: plus vite que la télévision}\label{tut:hf:prediction-and-sports-market-making:tut}

\begin{tutorial}
\textbf{Objectif.} Pricer un match en direct, le coter sur une bourse et mesurer ce que les courtsiders prennent à chaque \omterm{def:hf:prediction-and-sports-market-making:delay}{délai de pari}. \textbf{État final~:} les trois
figures et les chiffres du texte.
\begin{tutsteps}
\item \textbf{Le modèle de match}~: la probabilité d'une victoire à domicile à partir du score et du temps restant (\texttt{firm.sportsmm.home\_win}).
\item \textbf{La course à chaque but} et la perte du teneur de marché.
\omcode{code/firm/sportsmm/firm_sportsmm.py}{55}{75}{Le saut de la probabilité à un but~; l'ordre du courtsider face à l'annulation du teneur~; la perte sur la cotation périmée.}\label{lst:hf:prediction-and-sports-market-making:trade}
\item \textbf{Une saison} de 2\,000 matchs pour chaque délai et chaque plafond de mise (\texttt{hf\_sports}).
\end{tutsteps}
\textbf{Pour aller plus loin.} Suspendre le marché pendant le délai au lieu de maintenir les ordres~; ajouter un second flux de données, plus lent, pour une partie des
teneurs de marché~; ajouter les cartons rouges et les penalties comme événements.
\end{tutorial}

\section{Construction~: le teneur de marché sportif}\label{bld:hf:prediction-and-sports-market-making:sportsmm}

\begin{build}
\bfield{Rôle} Pricer des matchs en direct, coter sur une bourse avec un \omterm{def:hf:prediction-and-sports-market-making:delay}{délai de pari} et limiter les contreparties.
\bfield{Interface} \texttt{home\_win(lh, la, t\_left, score\_diff)}, \texttt{Match(lh, la, seed)},
\texttt{trade\_match(match, delay\_s, mm\_lat, fast\_lat, size, half\_spread, rec\_per\_min, rec\_stake, cap, lat\_sd)}, \texttt{season(n, delay\_s, seed)}, \texttt{odds\_quote(p, half\_spread)}.
Construit à partir de \texttt{firm.odds} du livre~3.
\bfield{Règles} Probabilités dans $[0,1]$~; mises exprimées en unités, chaque unité rapportant un~; latences lognormales autour de leurs médianes.
\bfield{Tests d'acceptation} \texttt{code/firm/sportsmm/tests/}~: domicile, nul et extérieur somment à un~; la fin d'un match résout la probabilité~; un but
la fait monter~; sans délai, chaque but est saisi~; avec dix secondes, aucun ne l'est~; un plafond proportionne la perte~; les cotes et l'overround de la cotation, calculés à la main.
\bfield{Pour aller plus loin} Suspension du marché~; cartons rouges~; flux lents et rapides entre teneurs de marché.
\end{build}

\begin{omsources}
\item K.~Croxson, J.~J.~Reade, Information and efficiency: goal arrival in soccer betting, \emph{Economic Journal} 124(575), 2014, 62--91.
\item Betfair developer support, ``Why do you have a delay on placing bets on a market that is in-play''.
\item ABC News, Australian Open illegal betting charges dropped, 6 mars 2014.
\item Le livre~3, chapitre~27, et ses sources pour l'encadré daté.
\end{omsources}

\section{Exercices}

\begin{exercise}[$\star$]\label{exo:hf:prediction-and-sports-market-making:1}
Un teneur de marché cote la victoire à domicile à 0.46 avec une demi-fourchette d'un point. Calculer les cotes décimales obtenues par un parieur de chaque côté, ainsi que l'overround.
\end{exercise}

\begin{exercise}[$\star$]\label{exo:hf:prediction-and-sports-market-making:2}
Un but fait passer la victoire à domicile de 0.46 à 0.70. Un courtsider prend 1\,000 unités à la cotation périmée de 0.47. Calculer la perte du teneur de marché.
\end{exercise}

\begin{exercise}[$\star$]\label{exo:hf:prediction-and-sports-market-making:3}
Le flux du teneur de marché signale les buts après trois secondes, et les courtsiders agissent après une demi-seconde. Déterminer le délai qui protège tout juste le teneur de marché, en
faisant abstraction de la variabilité.
\end{exercise}

\begin{exercise}[$\star\star$]\label{exo:hf:prediction-and-sports-market-making:4}
Expliquer pourquoi un délai de deux secondes laisse encore 72~\% des buts saisis.
\end{exercise}

\begin{exercise}[$\star\star$]\label{exo:hf:prediction-and-sports-market-making:5}
Expliquer pourquoi le \omterm{def:hf:prediction-and-sports-market-making:delay}{délai de pari} est un \omterm{def:hf:latency-arbitrage-and-its-defence:bump}{ralentisseur asymétrique}, et ce qui le rendrait symétrique.
\end{exercise}

\begin{exercise}[$\star\star$]\label{exo:hf:prediction-and-sports-market-making:6}
Expliquer ce que Croxson et Reade ont testé, et pourquoi les buts constituent un test propre.
\end{exercise}

\begin{exercise}[$\star\star\star$]\label{exo:hf:prediction-and-sports-market-making:7}
\emph{Programmation.} Diviser par deux la latence du flux de données du teneur de marché, la ramenant à 1.5~secondes. Quel délai rend désormais la perte négligeable~?
\end{exercise}

\begin{exercise}[$\star\star\star$]\label{exo:hf:prediction-and-sports-market-making:8}
\emph{Trouver la faille.} «~Notre modèle en direct est meilleur que celui du marché, donc nous pouvons coter plus serré que quiconque et gagner malgré tout.~»
\end{exercise}

\section{Problème~: plus rapide que la télévision}

\begin{problem}[{Problème du week-end --- plus rapide que la télévision}]\label{pb:hf:prediction-and-sports-market-making:1}
Un teneur de marché cote des marchés du vainqueur d'un match de football en direct sur une bourse assortie d'un \omterm{def:hf:prediction-and-sports-market-making:delay}{délai de pari}.

\medskip\noindent\textbf{Partie I --- Tarification.}
\begin{enumerate}
\item Écrire la probabilité de victoire à domicile, avec des buts suivant une loi de Poisson, et la donner au coup d'envoi, puis à la mi-temps avec un but d'avance.
\item Que trouvent Croxson et Reade~?
\item Convertir la cotation du teneur de marché en cotes et en overround.
\item Résumer l'encadré daté.
\end{enumerate}

\medskip\noindent\textbf{Partie II --- La course.}
\begin{enumerate}[resume]
\item Définir le \omterm{def:hf:prediction-and-sports-market-making:courtsiding}{courtsiding} et le \omterm{def:hf:prediction-and-sports-market-making:delay}{délai de pari}.
\item Décrire la course qui suit un but.
\item Donner la perte par but et la part des mises captées pour des délais de 0, 2, 3, 4 et 6 secondes.
\item Quand le marché paie-t-il le teneur de marché~?
\end{enumerate}

\medskip\noindent\textbf{Partie III --- Contreparties.}
\begin{enumerate}[resume]
\item Définir la \omterm{def:hf:prediction-and-sports-market-making:limiting}{limitation des mises}.
\item Que font les caps de 100 et de 50 unités~?
\item Que s'est-il passé dans l'affaire de \omterm{def:hf:prediction-and-sports-market-making:courtsiding}{courtsiding} de 2014~?
\item Pourquoi un teneur de marché opérant sur une bourse ne peut-il pas limiter lui-même les comptes~?
\end{enumerate}

\medskip\noindent\textbf{Partie IV --- Le verdict.}
\begin{enumerate}[resume]
\item Énoncer le \emph{résultat nommé}~: la perte du teneur de marché par but face aux courtsiders, en fonction du \omterm{def:hf:prediction-and-sports-market-making:delay}{délai de pari}, et le délai qui la rend
négligeable.
\item Qui paie pour le délai~?
\item Que changeraient des données officielles plus rapides~?
\item Quelle fiche de stratégie est la plus exposée à la réglementation~?
\item Comment mesurer la latence de votre flux par rapport à la vérité de terrain~?
\item Comment ce problème se compare-t-il à l'\omterm{def:hf:latency-arbitrage-and-its-defence:arb}{arbitrage de latence} sur actions (chapitre~9)~?
\item Que devrait faire une plateforme, en plus du délai~?
\item En une phrase~: que protège un \omterm{def:hf:prediction-and-sports-market-making:delay}{délai de pari}~?
\end{enumerate}
\end{problem}

\section{Questions d'entretien}

\begin{interviewq}[$\star$ \iqroles{trader}]\label{iq:hf:prediction-and-sports-market-making:1}
Une équipe mène 1--0 à dix minutes de la fin. Quelle devrait être à peu près sa probabilité de victoire si les deux équipes marquent 1.3~but par match~?
\end{interviewq}

\begin{interviewq}[$\star\star$ \iqroles{chercheur}]\label{iq:hf:prediction-and-sports-market-making:2}
Comment estimeriez-vous les taux de buts d'un match à partir des cotes d'avant-match~?
\end{interviewq}

\begin{interviewq}[$\star\star$ \iqroles{développeur}]\label{iq:hf:prediction-and-sports-market-making:3}
Concevez le système qui annule 2\,000 cotations réparties sur 300 matchs dans les 100~millisecondes suivant un but survenu dans l'un quelconque d'entre eux.
\end{interviewq}

\begin{interviewq}[$\star\star$ \iqroles{risque}]\label{iq:hf:prediction-and-sports-market-making:4}
Vos pertes en direct ont triplé dans une seule ligue. Qu'examinez-vous~?
\end{interviewq}

\begin{interviewq}[$\star\star$ \iqroles{trader}]\label{iq:hf:prediction-and-sports-market-making:5}
Pourquoi un bookmaker pourrait-il accueillir pendant un temps le compte d'un parieur aguerri avant de le limiter~?
\end{interviewq}

\begin{interviewq}[$\star\star\star$ \iqroles{chercheur}]\label{iq:hf:prediction-and-sports-market-making:6}
Avec des latences lognormales, dériver la probabilité qu'un ordre de courtsider, maintenu pendant le délai, arrive avant l'annulation du teneur de marché.
\end{interviewq}
